% !TEX root = ../main.tex
% =====================================================================
% CAPITOLO 1 — INTRODUZIONE
% =====================================================================

\chapter{Introduzione}
\label{chap:introduzione}

% =====================================================================
\section{Il contesto: dall'entusiasmo per gli LLM al problema dell'adozione enterprise}
\label{sec:contesto}
% =====================================================================

Il rilascio pubblico di ChatGPT nel novembre 2022 ha segnato l'inizio di una notevole adozione individuale dell'AI generativa: 
in meno di due mesi il servizio ha raggiunto 100~milioni di utenti, e nel giro di un anno strumenti basati su Large Language Model (LLM) 
sono entrati nella routine quotidiana di milioni di professionisti per la generazione di testi, 
l'analisi di documenti, la scrittura di codice e la risposta a domande più o meno complesse.

L'adozione \emph{enterprise}, tuttavia, ha seguito un percorso diverso, lento e incerto, 
principalmente a causa di queste criticità:

\paragraph{Costo.} Le licenze SaaS dei principali fornitori di AI generativa sono tariffate per utente. 
Per un'organizzazione con migliaia di dipendenti, il costo cumulativo di licenze nominali può raggiungere
decine e centinaia di migliaia di euro annui (anche alla luce dei prezzi che, con l'evolversi dei modelli, sembrano stare gradualmente crescendo), 
un investimento difficile da giustificare quando l'utilizzo effettivo dipende dall'approccio personale e rischia di essere limitato a una frazione della popolazione aziendale.
Inoltre il ritorno sull'investimento è ancora incerto e di difficile stima.

\paragraph{Sovranità del dato.} Le soluzioni pubbliche processano i dati dell'utente sui server del fornitore che possono trovarsi in data center non europei. 
Per le organizzazioni soggette al GDPR e, dal 2024, all'European AI Act~\cite{eu2024aiact}, la questione di \emph{dove} i dati vengono elaborati e \emph{chi} può accedervi è un vincolo 
rilevante.

\paragraph{Affidabilità.} I modelli generativi soffrono del fenomeno dell'\emph{allucinazione}: 
possono produrre affermazioni plausibili e ben formulate ma fattualmente false o errate~\cite{xu2024hallucination}. 
In un contesto aziendale, un'informazione errata su una procedura, un dato finanziario 
o un obbligo contrattuale può avere conseguenze concrete con effetti potenzialmente imprevedibili. 

\begin{figure}[htbp]
    \centering
    \includegraphics[height=0.85\textheight,keepaspectratio]{imbuto}
    \caption[Imbuto narrativo della tesi]{Imbuto narrativo: dal contesto generale dell'adozione enterprise dell'AI generativa ai contributi specifici di questo lavoro.}
    \label{fig:imbuto_narrativo}
\end{figure}

% =====================================================================
\section{Tre modi di ancorare un LLM alla realtà aziendale}
\label{sec:grounding-tre-modi}
% =====================================================================

La risposta al problema dell'allucinazione è il \emph{grounding}: l'ancoraggio della generazione del modello a fonti esterne verificabili, 
che forniscono al modello il contesto necessario per produrre risposte accurate. Il grounding opera a livello di inferenza (non modifica i pesi del modello) 
e si presenta in tre forme, ciascuna adeguata a un diverso tipo di fonte.

\paragraph{Web grounding.} Il modello accede ai risultati di un motore di ricerca per integrare informazioni attuali e pubbliche nella risposta.

\paragraph{Retrieval-Augmented Generation (RAG).} La domanda dell'utente è utilizzata per recuperare frammenti di documenti da un corpus indicizzato~\cite{lewis2020rag}; 
il modello genera la risposta vincolandosi a quel materiale.

\paragraph{Text-to-SQL.} La domanda è tradotta in una query SQL eseguita su un data warehouse; i risultati diventano il contesto per la generazione~\cite{yu2018spider,pourreza2023dinsql}.

I tre meccanismi sono descritti in dettaglio nel Capitolo~\ref{chap:background} (Sez.~\ref{sec:bg-grounding}).

La combinazione di queste tre tecniche di grounding \emph{all'interno di una stessa organizzazione} con l'ausilio di agenti specializzati per ciascun contesto 
e un meccanismo di coordinamento che li integri è una questione meno presente in letteratura, e costituisce il cuore di questa tesi.

% =====================================================================
\section{Il caso di studio}
\label{sec:caso-studio}
% =====================================================================

Il lavoro è stato svolto per intero nel contesto aziendale di
\href{https://www.coopalleanza3-0.it}{Coop Alleanza 3.0}, la 
\href{https://it.wikipedia.org/wiki/Coop_Alleanza_3.0}{più grande cooperativa di consumatori} del sistema Coop.
Lo stack IT dell'organizzazione comprende Google Cloud Platform come piattaforma cloud principale, 
Microsoft Entra ID come tenant identitario e SysAid come sistema di ticketing IT.

Per le ragioni citate (controllo dei costi, sovranità del dato, desiderio di avere una impronta etica che attenui il bias), è stato deciso dalla Cooperativa
di procedere con lo sviluppo di strumenti interni che consentano di sfruttare l'AI generativa alle proprie condizioni.

La domanda progettuale che ne deriva è: \emph{come si costruisce una piattaforma di AI generativa enterprise eterogenea, affidabile, economicamente sostenibile e sotto il pieno controllo dell'organizzazione?}

% =====================================================================
\section{Contributo e approccio}
\label{sec:contributi}
% =====================================================================

La tesi sostiene che la risposta è un'architettura ad \emph{agenti verticali specializzati}, 
ciascuno con la propria forma di grounding e il proprio dominio di competenza, 
\emph{coordinati da un orchestratore intelligente} 
che instrada dinamicamente le richieste dell'utente verso l'agente più appropriato.

I contributi specifici sono tre:

\begin{enumerate}
  \item \textbf{Progettazione e implementazione di una suite di agenti verticali su GCP.} 
  Quattro agenti specializzati: 
  Vera (generalista con Web grounding e RAG per-utente), CooPolicy (RAG su corpus documentale condiviso), 
  AllyCare (Text-to-SQL su dati NPS), SysAid (Text-to-SQL unificato con due fornitori di LLM) realizzati su infrastruttura cloud-native
  e dotati di autenticazione federata, catena di sicurezza e pipeline CI/CD. 
  Capitoli~\ref{chap:agenti-architettura} e~\ref{chap:agenti-implementazione}.

  \item \textbf{Progettazione, implementazione e valutazione dell'orchestrazione.} 
  Un orchestratore LLM-based che instrada le richieste sulla base di una stima di intento, complessità e confidenza, 
  con meccanismo di clarification quando l'ambiguità è alta. 
  La capacità generalista di Vera confluisce in Orion come opzione \texttt{self} (gestione delle richieste non specialistiche e fallback): 
  il router instrada quindi tra \texttt{self}, \texttt{sysaid}, \texttt{allycare} e \texttt{coopolicy}. 
  Capitoli~\ref{chap:orion-architettura} e~\ref{chap:orion-implementazione}.

  \item \textbf{Una metodologia di valutazione del routing.} 
  Basata su un gold standard configurabile e su un system prompt adattabile
  tramite interfaccia grafica per velocizzare le iterazioni e valutare velocemente
  le variazioni di performance.
  Capitolo~\ref{chap:orion-valutazione}.
\end{enumerate}

L'approccio si colloca nel filone della letteratura sui sistemi 
multi-agente basati su LLM~\cite{wu2023autogen,guo2024multiagent} 
e sul routing LLM~\cite{ong2024routellm}, con la specificità di un'implementazione 
completa e collocata in un contesto reale.

\begin{figure}[htbp]
    \centering
    \includegraphics[width=0.7\textwidth,height=0.85\textheight,keepaspectratio]{contributi}
    \caption[Sintesi visuale dei contributi della tesi]{Sintesi visuale dei tre contributi: suite di agenti verticali, orchestratore, metodologia di valutazione del routing.}
    \label{fig:contributi_sintesi}
\end{figure}

% =====================================================================
\section{Cosa questa tesi \emph{non} è}
\label{sec:non-e}
% =====================================================================

È utile delimitare il perimetro del lavoro svolto e delle scelte fatte:

\begin{itemize}
  \item Questa tesi \textbf{non propone una nuova architettura di modello}. 
  I modelli utilizzati (Gemini, Claude) 
  sono servizi gestiti erogati da fornitori commerciali; il lavoro si concentra sul modo 
  in cui sono utilizzati, orchestrati e valutati.

  \item Questa tesi \textbf{non propone un nuovo algoritmo di routing appreso}. 
  Il routing in Orion è LLM-based (un modello che classifica le richieste), 
  non un classificatore addestrato su dati di preferenza. 
  L'alternativa appresa è discussa come direzione futura.

  \item Questa tesi \textbf{non sostiene che l'architettura proposta sia universalmente la risposta giusta}. 
  Mostra che lo è \emph{in questo contesto}, alla luce degli specifici vincoli esposti.

\end{itemize}

% =====================================================================
\section{Struttura della tesi}
\label{sec:struttura}
% =====================================================================

La tesi è organizzata in otto capitoli.

Il \textbf{Capitolo~\ref{chap:background}} introduce i concetti necessari per apprezzare i contributi: 
architettura Transformer, capacità e limiti degli LLM, prompt engineering e output strutturato, 
grounding nelle tre forme (Web, RAG, Text-to-SQL), 
architetture multi-agente, LLM routing, stima dell'incertezza, sicurezza e compliance, 
identity federation e infrastruttura cloud.

Il \textbf{Capitolo~\ref{chap:agenti-architettura}} presenta i principi guida 
e le scelte architetturali della suite di agenti verticali: 
principi etici, governance, agnosticità, sovranità del dato, 
modularità verticale, automazione dei rilasci e controllo degli accessi.

Il \textbf{Capitolo~\ref{chap:agenti-implementazione}} descrive l'implementazione della suite: 
stack tecnologico, organizzazione multi-repository, infrastruttura Terraform, 
pipeline CI/CD, catena di sicurezza e dettagli implementativi dei quattro agenti.

Il \textbf{Capitolo~\ref{chap:orion-architettura}} presenta l'architettura dell'orchestratore: motivazione, 
routing step con output strutturato, logica di confidence e clarification, 
agent proxy, gestione dell'identità e della storia conversazionale.

Il \textbf{Capitolo~\ref{chap:orion-implementazione}} descrive l'implementazione dell'orchestratore: 
router agent, agent proxy, endpoint di chat con routing integrato, 
frontend con badge trasparente, admin panel con pagina di valutazione, 
infrastruttura Terraform e schema del database.

Il \textbf{Capitolo~\ref{chap:orion-valutazione}} presenta la valutazione sperimentale su tre piani: 
il routing dell'orchestratore (domande di ricerca, gold standard, metriche dichiarate a priori, 
risultati e iterazione del prompt, analisi degli errori e limiti), 
la qualità dei sub-agenti (execution accuracy per gli agenti Text-to-SQL e faithfulness per il RAG) 
e una valutazione economica a ordini di grandezza.

Il \textbf{Capitolo~\ref{chap:conclusioni}} sintetizza i contributi, discute la generalizzabilità del pattern e presenta le direzioni di sviluppo futuro.
